% Summary in English – a szerző írja.
Analysts in security operations centres receive large numbers of intrusion detection alerts, of which only a small share indicates a real attack. This overload leads to alert fatigue, in which genuine attacks may go unnoticed. The aim of this thesis is to design a system that groups and ranks network- and host-based intrusion detection alerts using metrics and information about the monitored environment, with the help of cooperating agents.

The literature review showed that each existing approach to alert grouping and prioritisation addresses only part of the problem. The strongest method that requires no training data, risk-based alert scoring, relies solely on the alert stream and ignores the importance, vulnerabilities and monitoring coverage of the affected host. Context-aware methods mostly filter alerts and may therefore discard real attacks when the asset inventory is incomplete, while solutions built on large language models are slow and not reproducible.

The proposed solution extends risk-based scoring with three context metrics that only change the order of alerts and never remove them. The metrics are combined in a weighted sum whose weights are set by pairwise comparison before any measurement takes place. The planned system consists of seven agents: the score and the ranking are produced by deterministic components, while a locally hosted language model only writes verified explanations. The effect will be measured with incrementally extended model variants on eight scenarios of a public alert dataset and in a dedicated laboratory environment, guarding against bias through pre-registration of settings and a control experiment.

The next semester is devoted to implementing the system, carrying out the measurements and evaluating the results.